\section{Round-twenty-three reconstruction: deterministic prediction, dominated observations, and persistently excited inference}
\label{sec:r23-c1}

The preceding revision attempted to dominate deterministic hidden transitions
by one probability measure.  This is impossible for an uncountable family of
Dirac kernels.  The active experiment keeps the deterministic transition as a
push-forward and imposes domination only on the observation channel.  Filter
continuity is proved by an integrated numerator identity that remains valid at
zero evidence and on strata of infinite reference measure.

\subsection{Typed hidden transition and observation strata}

Let $X$ be the prepared A3 history state or the prepared B2 collision-history
state, let $A$ be a compact action space, and let $\Theta\subset\mathbb R^d$
be compact.  For $(\vartheta,a)\in\Theta\times A$, the hidden transition is
\[
 M_\vartheta^a(x,dx')=\delta_{\Phi_\vartheta^a(x)}(dx'). \tag{C1.1}
\]
It is never assigned a density.  We assume the prepared flow
$(\vartheta,a,x)\mapsto\Phi_\vartheta^a(x)$ is continuous on every compact
Lyapunov shell supplied by A4 or B4.

The observation space is a disjoint union
\[
 Y=\bigsqcup_{s\in S}Y_s,\qquad
 \nu=\sum_{s\in S}2^{-|s|}\nu_s,                         \tag{C1.2}
\]
where $\nu_s$ is counting measure on lattice/atomic strata, surface measure on
contact strata, or Lebesgue measure on continuous aggregate strata.  The
observation kernel, and only this kernel, is dominated:
\[
 G_\vartheta^a(x,dy)=g_\vartheta^a(x,y)\nu(dy).           \tag{C1.3}
\]
On a prepared compact shell $K\subset X$ the map
$(\vartheta,a,x)\mapsto g_\vartheta^a(x,\cdot)$ is continuous in
$L^1(\nu)$.  Its first three parameter derivatives have integrable envelopes,
and the square-root derivative satisfies the uniform quadratic-mean bound
\[
 \int\left|\sqrt{g_{\vartheta+h}^a(x,y)}-\sqrt{g_\vartheta^a(x,y)}
 -\tfrac12h^\mathsf T\dot\ell_\vartheta^a(x,y)
       \sqrt{g_\vartheta^a(x,y)}\right|^2d\nu=o(|h|^2).  \tag{C1.4}
\]
No positive lower density is required on an infinite stratum.

\begin{proposition}[Model-derived observation regularity]
\label{prop:r23-c1-observation}
The Sinai observation blocks obtained from the A2 source-inserted local
operator and the hard-sphere blocks obtained from the B1 source-inserted
coefficient theorem satisfy (C1.2)--(C1.4) on the forward reachable compact
shells of A3--A4 and B2--B4.
\end{proposition}

\begin{proof}
For a Sinai block, keep the entrance history $x$ as an insertion in A2's
branchwise transfer operator.  The weighted branch derivative estimate is
uniform on an A4 Lyapunov shell, and the raw Fourier majorant is integrable.
Fourier inversion therefore gives an $L^1$ density on the continuous roof
stratum and counting densities on the three lattice strata, continuously in
$x$, with three source/parameter derivatives.  For a hard-sphere block, retain
the complete cut history from B2 at the left endpoint.  B1's deterministic
adapted block atlas gives a boundary-flat $W^{s,1}$ push-forward uniformly in
that history, and its exact-number coefficient ratio preserves the lattice and
surface atoms.  Differentiating the inserted pressure and applying Cauchy
bounds gives (C1.4).  These are conditional block calculations with the hidden
history held as data; an aggregate LLT is not re-labelled as a pointwise
kernel theorem.
\end{proof}

\subsection{Filter without common hidden-state domination}

For a belief $\pi\in\mathcal P(X)$ define the deterministic prediction
\[
 \pi^-=(\Phi_\vartheta^a)_\#\pi.                          \tag{C1.5}
\]
Its evidence density and unnormalised numerator are
\[
 r_\vartheta(\pi,a,y)=\int g_\vartheta^a(x,y)\pi^-(dx),
 \qquad
 N_\vartheta^{\varphi}(\pi,a,y)
 =\int\varphi(x)g_\vartheta^a(x,y)\pi^-(dx).             \tag{C1.6}
\]
When $r>0$, the posterior $B_\vartheta(\pi,a,y)$ is defined by
$B(\varphi)=N^\varphi/r$.  When $r=0$, set it equal to an arbitrary fixed
belief $\pi_\dagger$; this convention is never observed because its weight in
the belief kernel is zero.  The belief transition is
\[
 Q_\vartheta^aF(\pi)
 =\int_YF(B_\vartheta(\pi,a,y))
       r_\vartheta(\pi,a,y)\nu(dy).                      \tag{C1.7}
\]

\begin{theorem}[Integrated Feller filter]
\label{thm:r23-c1-filter}
On every prepared compact belief shell, $(\vartheta,a,\pi)\mapsto
Q_\vartheta^aF(\pi)$ is continuous for every bounded continuous
$F:\mathcal P(X)\to\mathbb R$.  The conclusion does not require a common
dominating measure for (C1.1), a lower observation density, or finite
$\nu(Y)$.
\end{theorem}

\begin{proof}
Let $(\vartheta_n,a_n,\pi_n)\to(\vartheta,a,\pi)$.  Continuity of the
deterministic flow gives $\pi_n^-\Rightarrow\pi^-$.  The $L^1(\nu)$
continuity and integrable envelope in (C1.3) imply
$r_n\to r$ in $L^1(\nu)$.  For every bounded Lipschitz $\varphi$,
$N_n^\varphi\to N^\varphi$ in $L^1(\nu)$.  With ratios set to zero when their
denominator vanishes,
\[
 r_n\left|{N_n^\varphi\over r_n}-{N^\varphi\over r}\right|
 \le |N_n^\varphi-N^\varphi|+\|\varphi\|_\infty|r_n-r|. \tag{C1.8}
\]
Choose a countable bounded-Lipschitz family metrising weak convergence of
beliefs, sum (C1.8) with geometric weights, and obtain
\[
 \int r_n(y)d_{\rm BL}(B_n(y),B(y))d\nu(y)\longrightarrow0.             \tag{C1.9}
\]
Uniform continuity of $F$ on the compact reachable belief shell and
$\|r_n-r\|_1\to0$ now prove (C1.7).  This weighted-ratio identity controls
small evidence directly; no estimate of the form $\delta\nu(Y)$ occurs.
\end{proof}

\subsection{Risk-sensitive dynamic programming}

For a bounded continuous one-step reward $c(\pi,a,y)$ and terminal reward
$\Psi$, define the finite-horizon exponential value in unnormalised form.  The
Markov state is the belief $\pi$ together with accumulated log normaliser; the
one-step operator is
\[
 (\mathcal T_\vartheta V)(\pi)
 =\inf_{a\in A}\int
 \exp\{c(\pi,a,y)\}V(B_\vartheta(\pi,a,y))
 r_\vartheta(\pi,a,y)d\nu(y).                            \tag{C1.10}
\]
Theorem~\ref{thm:r23-c1-filter}, compactness of $A$, and measurable selection
give continuous values and Borel optimal selectors.  Zero evidence makes zero
contribution and does not create a projective state.

\subsection{Finite coordinates in the convex belief space}

Let $\mathcal K\subset\mathcal P(X)$ be a compact reachable belief set.  Pick
an $\epsilon$-net $\pi_1,\ldots,\pi_J$ in bounded-Lipschitz distance and a
finite family of bounded Lipschitz functions $\varphi_1,\ldots,\varphi_m$ that
separates the net with margin.  Put
\[
 T_m(\pi)=(\pi\varphi_1,\ldots,\pi\varphi_m).             \tag{C1.11}
\]
Choose a partition of unity $\psi_i$ on a neighbourhood of $T_m(\mathcal K)$,
subordinate to coordinate balls around $T_m(\pi_i)$, and define
\[
 R_m(z)=\sum_{i=1}^J\psi_i(z)\pi_i.                       \tag{C1.12}
\]
Because $\mathcal P(X)$ is convex, $R_m(z)$ is always a belief, even if the
reachable set $\mathcal K$ is nonconvex.

\begin{proposition}[Finite-coordinate value approximation]
\label{prop:r23-c1-finite}
The construction can be chosen so that
\[
 \sup_{\pi\in\mathcal K}d_{\rm BL}(R_mT_m\pi,\pi)
 \le3\epsilon.                                           \tag{C1.13}
\]
Consequently every uniformly continuous finite-horizon value is uniformly
approximated by the finite-coordinate dynamic program obtained by applying
$R_m$ after each coordinate update.
\end{proposition}

\begin{proof}
If $\psi_i(T_m\pi)>0$, separation of the net implies
$d_{\rm BL}(\pi_i,\pi)\le3\epsilon$.  Convexity of the metric under mixtures
gives (C1.13).  Iterating the Feller modulus for finitely many steps controls
the value error.  No retraction into an arbitrary nonconvex compact set is
claimed.
\end{proof}

\subsection{Adaptive LAN under persistent excitation}

Let $\alpha=(a_k)$ be a predictable policy.  The likelihood increment is the
evidence density from (C1.6), evaluated at the parameter-dependent filter:
\[
 \ell_k(\vartheta)=
 \log r_\vartheta(\pi_{k-1}^{\vartheta},a_k,Y_k).         \tag{C1.14}
\]
Differentiability of the filter follows by differentiating the unnormalised
numerator before normalisation and using (C1.4) together with the integrated
identity (C1.8).  The score
$s_k=\partial_\vartheta\ell_k(\vartheta_0)$ is a martingale difference for the
observation filtration.

The policy class $\mathfrak A_{\rm pe}$ is defined by the explicit persistent
excitation and identifiability conditions
\[
 cI_d\le {1\over n}\sum_{k=1}^n
 E_{\vartheta_0}[s_ks_k^\mathsf T\mid\mathcal F_{k-1}]
 \le CI_d+o_P(1),                                        \tag{C1.15}
\]
uniformly in $\alpha$, and by a uniform Kullback separation outside every
fixed neighbourhood of $\vartheta_0$.  Policies which always choose an
uninformative action are not in $\mathfrak A_{\rm pe}$; compactness alone is
never said to imply (C1.15).

The derivative envelope in (C1.4) gives a uniform conditional Lindeberg bound
and a third-order remainder
\[
 \sup_{|h|\le R}\left|
 \sum_{k=1}^n\bigl[\ell_k(\vartheta_0+h/\sqrt n)
 -\ell_k(\vartheta_0)igr]
 -h^\mathsf T\Delta_n+\tfrac12h^\mathsf TI^\alpha h
 \right|\xrightarrow{P}0,                               \tag{C1.16}
\]
where $\Delta_n=n^{-1/2}\sum_ks_k$ and
$I^\alpha$ is the limit of (C1.15).

\begin{theorem}[Uniform adaptive LAN and Bernstein--von Mises]
\label{thm:r23-c1-bvm}
Uniformly over $\alpha\in\mathfrak A_{\rm pe}$,
$\Delta_n\Rightarrow N(0,I^\alpha)$ and (C1.16) holds.  If the prior has a
continuous positive density at $\vartheta_0$, then the posterior law of
$\sqrt n(\vartheta-\vartheta_0)$ converges in total variation to
\[
 N\bigl((I^\alpha)^{-1}\Delta_n,(I^\alpha)^{-1}\bigr).    \tag{C1.17}
\]
\end{theorem}

\begin{proof}
The martingale central limit theorem applied to (C1.15) and the conditional
Lindeberg estimate gives the score limit.  Taylor expansion of the recursively
differentiated evidences gives (C1.16); the third derivatives are summable by
the common envelope.  Uniform Kullback separation constructs exponentially
consistent tests on a finite cover of every fixed-distance alternative, while
(C1.16) gives the local denominator lower bound.  Posterior contraction to an
$n^{-1/2}$ neighbourhood follows.  On that neighbourhood, dominated
convergence in the LAN likelihood ratio and Gaussian tail control yield
(C1.17), following the standard LAN-to-BvM argument
\cite{vanDerVaart1998,LeCamYang2000}.  Every uniform step uses either
(C1.4), (C1.15), or the stated separation condition.
\end{proof}

\subsection{Round-Twenty-Two regression}

Distinct hidden states may have mutually singular transitions
$\delta_{\Phi(x)}$; this is permitted because (C1.1) is never dominated.
Only observations use $\nu$.  Aggregate local theorems enter through the
conditional inserted calculation in Proposition
\ref{prop:r23-c1-observation}.  Infinite observation strata and zero evidence
are handled by (C1.8), finite coordinates reconstruct into the convex belief
space rather than a nonconvex reachable set, and the information lower bound
is an explicit persistent-excitation condition verified policy by policy.
